% C.3bi finalized after full owner, root, frobenius and cycles proof/source audits.
\paragraph{C.3bi (Proved): a higher integral kernel in characteristic five.}
For $k=\overline{\mathbb F}_5$ there is an
excellent normal strictly henselian local $k$-domain $R$ of dimension five,
with residue field $k$, together with $0\ne f\in\mathfrak m_R$, such that
\[
 X=\operatorname{Spec}R\setminus\{\mathfrak m_R\},\qquad
 Y=X\cap V(f)
\]
are regular and the restriction
$\mathrm{CH}^2(X)_{\rm alg}\to\mathrm{CH}^2(Y)_{\rm alg}$ has a nonzero
class of order five in its kernel. Algebraic equivalence here uses actual
families on fixed products with smooth proper curves. In particular the
conclusion concerns integral coefficients; the exhibited class vanishes
after tensoring with $\mathbb Q$.

\begin{proof}
\textbf{1. A cubic surface and its divisor cup product.}
Let $W\subset\mathbb P^3_k$ be the Fermat cubic surface, let
$A=\mathcal O_W(1)$, and let $D\subset W$ be the line
\[
 w_0+w_1=w_2+w_3=0.
\]
The cubic is smooth since $3\ne0$ in $k$. Set
$u=w_0+w_1$, $v=w_2+w_3$, $s=w_1$, $t=w_3$. Its equation is
$uQ_0+vQ_1=0$, with $Q_0|_D=3s^2$ and $Q_1|_D=3t^2$.
The map
\[
 \mu_D:H^0(W,A)\longrightarrow H^1(W,\Omega^1_{W/k}\otimes A),
 \qquad a\longmapsto a\,c_1(\mathcal O_W(D)),
 \tag{C.3bi.1}
\]
has rank two. We give the precise calculation so that its integral
normalization is visible. The conormal and Euler sequences identify its
target with $H^2(W,\mathcal O_W(-2A))=H^0(W,A)^\vee$.
Indeed $H^1(W,\mathcal O_W(nA))=0$ for every integer $n$,
$H^2(W,\mathcal O_W(nA))=0$ for $n\ge0$, and
$H^0(\mathcal O_W)^4\to H^0(A)$ is an isomorphism. The normal
sequence of $D\subset W\subset\mathbb P^3$, twisted by $-A$, is
\[
 0\longrightarrow\mathcal O_D(-2)\longrightarrow\mathcal O_D^2
 \xrightarrow{(3s^2,3t^2)}\mathcal O_D(2)\longrightarrow0.
\]
Thus the obstruction to lifting $D$ in a deformation given by a quadratic
$q$ is $q|_D$ modulo $\langle s^2,t^2\rangle$.
The divisor boundary
\[
 H^1(D,\mathcal O_D(-2))\longrightarrow H^2(W,\mathcal O_W(-A))
\]
is an isomorphism: in
$0\to\mathcal O_W(-A)\to\mathcal O_W(D-A)\to\mathcal O_D(-2)\to0$,
$H^1$ and $H^2$ of the middle term vanish by
$D^2=-1$, $A\cdot D=1$, surface Riemann--Roch and Serre duality.
Lifting local defining equations of $D$ identifies this boundary with
contraction of the deformation cocycle with $d\log\mathcal O_W(D)$.
Consequently the Serre-dual bilinear form of (C.3bi.1), up to one common
nonzero scalar, is
\[
 (a,b)\longmapsto\lambda_D((ab)|_D),\qquad
 \lambda_D(s^2)=\lambda_D(t^2)=0,\quad\lambda_D(st)=1.
 \tag{C.3bi.2}
\]
It has rank two in the basis $(u,v,s,t)$. This calculation has no
factorial denominator and in particular is nonzero in characteristic five.

\textbf{2. The specified local deformation and its principal cut.}
Use the ordinary Fermat quartic and the quartic deformation from C.3bf:
\[
 C: F=z_0^4+z_1^4+z_2^4=0,\qquad
 H=z_0^2z_1^2+z_0^2z_2^2+z_1^2z_2^2.
\]
Write $K_C=\mathcal O_C(1)$ and $J=\operatorname{Pic}^0(C)$.
Choose $\lambda\in\mathbb F_{125}\subset k$ satisfying
$\lambda^3+\lambda+1=0$. This polynomial has no root in $\mathbb F_5$,
and reduction gives $\lambda^{12}=2-\lambda\notin\{1,3,4\}$.
Put $S=W\times C$ and choose the actual sections
\[
 a_0=z_0w_0+z_1w_1,\qquad a_1=z_0w_2+\lambda z_2w_3
 \quad\hbox{of }A\boxtimes K_C.
\]
Set
\[
 V=\mathbb P^1\times W\times C,\qquad
 L=\mathcal O_{\mathbb P^1}(1)\boxtimes A\boxtimes K_C,
 \qquad\ell=x_0a_0+x_1a_1.
 \tag{C.3bi.3}
\]
Here is an all-geometric-point verification of the smooth base locus.
The two rows $(z_0,z_1,0,0)$ and $(0,0,z_0,\lambda z_2)$ are
independent everywhere on $C$. Their common projective kernel is
$C\times\mathbb P^1$, parametrized by
\[
 (w_0,w_1,w_2,w_3)=(z_1\sigma,-z_0\sigma,
                      \lambda z_2\theta,-z_0\theta).
\]
The remaining cubic equation is
\[
 A_0\sigma^3+B_0\theta^3=0,\qquad
 A_0=z_1^3-z_0^3,\quad B_0=\lambda^3z_2^3-z_0^3.
\]
Each of $A_0$ and $B_0$ has twelve simple zeros on $C$, and their
zero divisors are disjoint. For $A_0$, set $z_1/z_0=\zeta$ with
$\zeta^3=1$; the fourth-power equation for $z_2/z_0$ has nonzero
right side $-(1+\zeta)$. For $B_0$, the corresponding right side
could vanish only if $\lambda^{12}=4$. A common zero would give
$1+\zeta+\lambda^{-4}\xi=0$ with $\zeta^3=\xi^3=1$, implying
$\lambda^{12}=3$ if $\zeta=1$ and $\lambda^{12}=1$ otherwise.
All are excluded. Away from these simple disjoint branch divisors,
the fibre derivative of $A_0\sigma^3+B_0\theta^3$ is nonzero;
at a branch point its base derivative is nonzero. Thus
$B=(a_0=a_1=0)\subset S$ is a smooth codimension-two complete
intersection, and $da_0,da_1$ are independent there.

Write $\ell=g_0z_0+g_1z_1+g_2z_2$, where
\[
 g_0=x_0w_0+x_1w_2,\qquad g_1=x_0w_1,\qquad
 g_2=\lambda x_1w_3.
 \tag{C.3bi.4}
\]
These sections are independent: the coefficient minor on
$x_0w_0,x_0w_1,x_1w_3$ is $\lambda$.
On $D$ their restrictions are $-x_0s-x_1t$, $x_0s$, and
$\lambda x_1t$. By (C.3bi.2), their cup products with $c_1(D)$
span a space of rank two. This fixes every coefficient over
$\mathbb F_{125}$; no generic or transcendental coefficient is used.

In the total space of
$\mathcal O_{\mathbb P^1\times W\times\mathbb P^2}(-1,-1,-1)$,
with fibre coordinate $\tau$, take the hypersurface
\[
 Z_0:\quad F+\tau\ell H=0.
 \tag{C.3bi.5}
\]
All terms have the same multidegree. Contract the zero section by the
Segre coordinates $\tau x_iw_jz_h$, take the reduced affine image,
normalize it, localize at the unique point above the vertex, and take its
strict henselization. This defines $R$ canonically from the chosen
coefficients. Let $\mathcal Z\to\operatorname{Spec}R$ be the induced
proper morphism. The verification from C.3bf applies with the additional
smooth factor $W$: the differential of $F$ is nonzero along $C$, so
$Z_0$ is regular near $V$; its exceptional fibre is the connected smooth
$V$, and normalization gives one point above the vertex. The contraction
is an isomorphism off that fibre. The singular locus away from the
regular neighbourhood of $V$ has proper image avoiding the vertex, and
therefore disappears after this strict localization. Hence $\mathcal Z$
is regular, of dimension five, and is isomorphic to $X$ off $V$.
The Segre generators have no common zero on the projective base, giving
\[
 \mathfrak m_R\mathcal O_{\mathcal Z}=\mathcal I_V,
 \qquad\mathcal I_V/\mathcal I_V^2=L.
 \tag{C.3bi.6}
\]
In the Segre coordinates $v_{a,i,j}=\tau x_aw_iz_j$, the exact
linear combination
\[
 f=v_{0,0,0}+v_{0,1,1}+v_{1,2,0}+\lambda v_{1,3,2}=\tau\ell
\]
is a
nonzero element of $\mathfrak m_R$. Its strict transform is the split
line bundle over
\[
 H_0=(\ell=0)\subset V.
\]
Indeed (C.3bi.5) becomes exactly $F=0$ there, at every order. The
incidence variety $H_0$ is the blow-up of $S$ in the smooth complete
intersection $(a_0,a_1)$, so is smooth. Its line-bundle puncture is $Y$;
thus $X$ and $Y$ are regular. The projection $\mathcal Z\to W$
exists on this actual nonsplit deformation. It extends the divisor $D$,
which will be essential below.

\textbf{3. An actual ordinary Picard kernel with a nonzero cup product.}
The hypersurface sequence for the cubic, Serre duality on $C$, and
K\"unneth give
\[
 H^1(V,L)=k^8,\quad H^2(V,L)=0,\qquad
 H^1(V,L^n)=H^2(V,L^n)=0\quad(n\ge2).
 \tag{C.3bi.7}
\]
Put $E=V_2\subset\mathcal Z$. As in C.3bf, its connected Picard group
is an extension
\[
 0\longrightarrow H^1(V,L)\longrightarrow G
 \longrightarrow J\longrightarrow0.
 \tag{C.3bi.8}
\]
The deformation class is the product of $\ell$ with the same quartic
normal deformation as in C.3bf. The nonsingular Serre-dual cup form of
that quartic identifies the image of its extension-class map with
\[
 U=\ell H^1(V,\mathcal O_V)\subset H^1(V,L),\qquad\dim U=3.
\]
This is the transition-unit/Atiyah calculation of
\citet[Theorem 2.6 and formula (2.7.2)]{D_Illusie1972Torsors}, with
all the auxiliary factors carried along. The first-neighbourhood
extension has exactly the same rank-three class as the indicated
three-dimensional subspace $U$.

The quartic has Hasse--Witt matrix $2I_3$ in characteristic five, as
certified in C.3bf. Thus $v:J^{(5)}\to J$ is an \'etale
Verschiebung with kernel $K=(\mathbb Z/5)^3$, and
$v^*:H^1(J,\mathcal O_J)\to H^1(J^{(5)},\mathcal O_{J^{(5)}})$
is zero. The integral, all-characteristic identification
$\operatorname{Ext}^1(J,\mathbb G_a)=H^1(J,\mathcal O_J)$
\citep[Section 3.2.3, equation (3.8)]{D_BrionIsogenyII}
gives a group lift $b:J^{(5)}\to G$. The exact Hom--Ext sequence
for $v$ identifies the coefficient map of $b|_K$ with the extension
class. It is injective on $K$ and its $k$-linear span is $U$.
By the explicit rank-two cup calculation of Step 2, there is $\eta\in K$ such that
\[
 a=b(\eta)\in U,\qquad
 a\smile c_1(D)\ne0
 \quad\hbox{in }H^2(V,L\otimes\Omega^1_{V/k}).
 \tag{C.3bi.9}
\]

We spell out why this point is represented by an actual ordinary
class. Choose a smooth connected proper curve $P\subset J^{(5)}$
through $0$ and $\eta$. Pull the normalized Poincar\'e bundle to
$V\times P$ through $v|_P$. The only first-lift obstruction is the
mixed Atiyah class, and its $H^1(J,\mathcal O_J)$ factor is killed by
$v^*$. Therefore it lifts to a bundle on $E\times P$. The difference
between its classifying map and $b|_P$ is a morphism from the proper
connected curve $P$ to the vector group $H^1(V,L)$; it is constant.
Tensoring by the inverse of the fibre at $0$ normalizes this constant
away. This gives exactly the classes $0$ and $a$ at the two endpoints.
For $n\ge2$,
\[
 H^2(V\times P,L^n)=0
\]
by (C.3bi.7), K\"unneth and $H^2(P,\mathcal O_P)=0$. The actual
bundle therefore lifts successively to every exceptional thickening.
Proper formal existence on $\mathcal Z_{\widehat R}\times P$
algebraizes the compatible family. Descend the bundle, its inverse and
its first-jet isomorphism to one smooth $R$-algebra in a Popescu
presentation of $\widehat R$; henselian smooth lifting modulo
$\mathfrak m_R^2$ gives an actual bundle on $\mathcal Z\times P$
with those same first jets. These are exactly the finite-data descent
and lifting steps proved in C.3bf. Endpoint formal classes are uniquely
determined by their first jets because of (C.3bi.7). Elkik's punctured
Picard comparison then proves that their restrictions to $X$ give the
ordinary class represented by $a$.

On the target, the blow-up description gives
$\operatorname{Pic}^0(H_0)=J$. The sequence for $H_0\subset V$ gives
\[
 H^1(H_0,L|_{H_0})=H^1(V,L)/U,
 \qquad H^1(H_0,L^n|_{H_0})=H^2(H_0,L^n|_{H_0})=0\ (n\ge2).
\]
For the stated $H^2$-vanishing use also $H^3(V,L^{n-1})=0$, which
follows from the cubic and curve cohomology above. The target
neighbourhood is split, so its connected formal Picard group is
$J\times(H^1(V,L)/U)$. The proper map induced by $b$ is $(v,0)$.
Consequently $a$ restricts to zero on the target at the first jet, then
at every order, and finally on the actual $Y$ by Elkik. The same argument applied to $5a=0$ proves first that the fifth
tensor power is trivial on $E$, and then on every exceptional
thickening by (C.3bi.7). It is therefore trivial on the completed
puncture. Elkik's injective comparison proves that its restriction
to the actual $X$ is trivial. Thus the actual class, not merely its
first jet, is killed by five. Let $\mathcal A$ be the endpoint bundle on
$\mathcal Z$, normalized so that
\[
 \mathcal A|_V=\mathcal O_V,\qquad
 \mathcal A|_E\text{ has relative first-jet class }a.
 \tag{C.3bi.10}
\]

\textbf{4. The relative $K_1$ boundary is zero.}
We use ordinary nonconnective algebraic $K$-theory and write
$K(E,V)=\operatorname{fib}(K(E)\to K(V))$. For the smooth proper
curve $C$, Quillen localization \citep[Section 7, Proposition 3.2 and (3.3),
p.~127]{rootQuillen1973}
gives the exact segment
\[
 \bigoplus_{c\in C(k)}k^*\longrightarrow K_1(C)
 \longrightarrow k(C)^*\xrightarrow{\operatorname{div}}
 \bigoplus_{c\in C(k)}\mathbb Z.
 \tag{C.3bi.11}
\]
The last kernel is $k^*$; both $k^*$ and the quotient of the first
direct sum by the preceding image are divisible groups. Thus $K_1(C)$
is divisible. A smooth cubic surface over an algebraically closed field
is a blow-up of $\mathbb P^2$ in six points
\citep[Chapter 4, Proposition 2.4, author p.~178]{D_HuybrechtsCubics};
applying the integral
projective-bundle and regular-blow-up formulas after product with $C$
gives
\[
 K_1(V)\simeq K_1(C)^{18}.
 \tag{C.3bi.12}
\]
The formulas used here are spectrum decompositions, not rational Chern
characters; for example see \citet[Corollaries 3.6 and 4.4, author
pp.~12 and 16]{D_Khan2020Blowup}. Hence $K_1(V)$ is five-divisible.

On each affine of a finite affine cover of the separated noetherian
scheme $E$, the ideal defining $V$ is square zero and the ring has
characteristic five. By \citet[Theorem A, p.~29]{rootGeisserHesselholt2011},
each relative $K$-group on such an affine, in every required degree,
is killed by a bounded power of five. Iterating the finite
Mayer--Vietoris sequences of \citet[Theorem 8.1]{D_ThomasonTrobaugh1990}
shows that $K_0(E,V)$ itself is killed by a bounded power of five.
The boundary from the five-divisible group $K_1(V)$ must therefore
be zero. In particular
\[
 K_0(E,V)\lhook\joinrel\longrightarrow K_0(E)
 \tag{C.3bi.13}
\]
is injective. This step is necessary: a relative infinitesimal class
alone would not establish a nonzero absolute $K_0$ class.

\textbf{5. The nonsplit first jet detects the product in relative $K_0$.}
For the square-zero pair $(E,V)$, write $I=L$ for its ideal.
The relative $K$-sheaves in degrees zero, one and two are
\[
 \mathcal K_0(E,V)=0,\qquad
 \mathcal K_1(E,V)=I,\qquad
 \mathcal K_2(E,V)=I\otimes\Omega^1_{V/k}.
 \tag{C.3bi.14}
\]
Here the last assertion is integral because $2$ is invertible.
Locally the square-zero algebra splits since $V$ is smooth; the
dual-number theorem of \citet[unnumbered theorem and corollary,
p.~A1204]{frVanDerKallen1971} gives the last identification. We choose
its sign so that
\[
 \{1+\tau c,u\}\longmapsto\tau c\,d\log u.
\]
This is the negative of the Dennis--Stein normalization in
\citet[Chapter III, Exercise 5.14]{frWeibelKbookIII}. Changing a local
splitting changes the chosen lift of $u$ by an element of $I$.
Its extra symbol has both entries infinitesimal and is zero when
$2$ is invertible. Changing the frame of $I$ just rescales its
coefficient. Thus (C.3bi.14) glues canonically on this \emph{nonsplit}
neighbourhood. Since $k$ is perfect, the absolute differentials over
$\mathbb Z$ in the cited theorem equal the differentials over $k$.

The relative descent spectral sequence, obtained by taking the fibre
in \citet[Theorem 10.3]{D_ThomasonTrobaugh1990}, is
\[
 E_2^{q,t}=H^q(V,\mathcal K_{-t}(E,V))
 \ \Longrightarrow\ K_{-q-t}(E,V).
 \tag{C.3bi.15}
\]
The underlying space is the four-dimensional noetherian $V$, so
its cohomological filtration has only the five steps $q=0,\ldots,4$
and is strongly convergent. A nonzero surviving associated-graded
class therefore detects a nonzero element; no infinite filtration
or inverse-limit assertion is involved.
It is a module over the absolute spectral sequence under tensor product
\citep[Section 3.15]{D_ThomasonTrobaugh1990}. The relative class of
$[\mathcal A|_E]-1$ has first term $a\in H^1(V,I)$.
Multiplying it by $[\mathcal O(D)|_E]-1$ gives a class $\widetilde x$
in relative filtration two whose term at $E_2^{2,-2}$ is exactly
\[
 a\smile d\log\mathcal O(D)\ne0.
 \tag{C.3bi.16}
\]
This follows directly as well by multiplying the transition functions
$1+\tau a_{ij}$ and the divisor transition functions: their relative
Steinberg symbols give the indicated cup cocycle.
There is only one possible incoming differential to this position,
namely $d_2:H^0(V,I)\to H^2(V,I\otimes\Omega^1_{V/k})$.
Every $u\in H^0(V,I)$ is represented by the actual global relative
unit $1+u$, and hence is permanent; this differential is zero.
There are no other incoming differentials since their first index
would be negative. The displayed product is a product of permanent
classes, so its outgoing differentials vanish. Therefore (C.3bi.16)
is its nonzero associated-graded class in $K_0(E,V)$.
By (C.3bi.13), its image
\[
 x_E=([\mathcal A|_E]-1)([\mathcal O(D)|_E]-1)
 \tag{C.3bi.17}
\]
is nonzero in $K_0(E)$.

More precisely, put $x=([\mathcal A]-1)([\mathcal O(D)]-1)$ in
$K_0(\mathcal Z)$. For each integer $j\ge0$, the pullback of
$(\psi^2)^j x$ to $E$ is
\[
 ([\mathcal A|_E]^{2^j}-1)
 ([\mathcal O(D)|_E]^{2^j}-1).
\]
Its unique relative lift, by (C.3bi.13), has filtration-two term
$4^j a\smile c_1(D)$. This uses only the ring and line-bundle
properties of ordinary Adams operations and their pullback
compatibility; it asserts no additional relative Adams theorem.

\textbf{6. The exceptional image cannot remove the Chow class.}
Let
\[
 \beta=c_1(\mathcal A)c_1(\mathcal O(D))\in\mathrm{CH}^2(\mathcal Z).
\]
Its restriction to $X$ is ordinary algebraically trivial because the
actual family of Step 3 can be intersected with the fixed divisor
pulled back from $W$. Its restriction to $Y$ is zero and it is
annihilated by five on $X$, by Step 3.
Suppose for contradiction that $\beta|_X=0$. Localization on the
regular pair $(\mathcal Z,V)$ gives $\beta=i_*\gamma$ with
$\gamma\in\operatorname{Pic}(V)$. Its restriction to $V$ is zero
because $\mathcal A|_V$ is trivial. The self-intersection formula
then gives $L\gamma=0$ in $\mathrm{CH}^2(V)$.

Multiplication by $L$ on $\operatorname{Pic}(V)$ is injective
integrally. Indeed, with $h=c_1(\mathcal O_{\mathbb P^1}(1))$ and
$c=c_1(A\boxtimes K_C)$ on $S$, write
$\gamma=\gamma_0+b h$. The projective-bundle formula yields
\[
 (h+c)(\gamma_0+bh)=c\gamma_0+h(\gamma_0+bc).
\]
Its vanishing gives $\gamma_0=-bc$ and $bc^2=0$. Intersecting the
last equality with $c$ and taking degree gives $b\deg(c^3)=0$ in
$\mathbb Z$. Since $c$ is ample, its degree is positive, so $b=0$
and $\gamma_0=0$. This argument retains any torsion in
$\operatorname{Pic}^0(C)$. Thus $\gamma=0$ and $\beta=0$ already
in $\mathrm{CH}^2(\mathcal Z)$.

\textbf{7. Integral Adams separation completes the contradiction.}
Let $F^qK_0(\mathcal Z)$ be the coniveau filtration. The integral
cycle map takes the product $\beta$ to $x$ modulo $F^3$; compatibility
with two divisor intersections also follows directly from the Cartier
exact sequences. Hence the conclusion of Step 6 implies $x\in F^3$.
The scheme $\mathcal Z$ is separated regular noetherian of dimension
five. Consequently
\[
 (\psi^2-2^q)F^q\subset F^{q+1},\qquad F^6=0,
\]
by \citet[Lemma 3.10 and Proposition 3.4]{D_IwasaKai2021}.
This is their \emph{absolute} coniveau lemma; their relative Chow
with modulus theorem and its affine-neighbourhood hypothesis are not
being used. Their Theorem 4.2 also supplies the integral cycle map
(and in codimension two it is an isomorphism onto the associated
graded). It follows that
\[
 P(\psi^2)x=0,\qquad P(T)=(T-8)(T-16)(T-32).
 \tag{C.3bi.18}
\]
Pull back to $E$. By Step 5 the unique relative lift of the left
side has filtration-two term
\[
 P(4)\,a\smile c_1(D)=-1344\,a\smile c_1(D)
                 =a\smile c_1(D)\ne0
 \quad\text{in characteristic five}.
\]
This contradicts (C.3bi.13) and (C.3bi.18). Therefore
$\beta|_X\ne0$. It is killed by five, so has exact order five.
All constructions of algebraic triviality used actual fixed-product
families, and the target vanishing was checked on the actual puncture.
This proves the claimed integral higher restriction failure.
\end{proof}

\noindent\emph{Scope.} This result concerns characteristic five,
codimension two and dimension five. No rational failure or
all-dimensional extension follows from the displayed exponent-five
class. The exact finite certificates and complete independent proof reviews
are retained with this statement. They certify the scopes stated in
their separate logs; the finite computations do not replace the
global proof.
